
\documentclass[10pt,a4paper]{article}
\usepackage[a4paper,margin=1.65cm]{geometry}
\usepackage[utf8]{inputenc}
\usepackage[T1]{fontenc}
\usepackage[english]{babel}
\usepackage{lmodern}
\usepackage{microtype}
\usepackage{xcolor}
\usepackage{listings}
\usepackage{booktabs}
\usepackage{longtable}
\usepackage{array}
\usepackage{enumitem}
\usepackage{hyperref}
\usepackage{fancyhdr}
\usepackage{titlesec}
\usepackage{amsmath}
\definecolor{codebg}{RGB}{248,248,248}
\definecolor{codeframe}{RGB}{215,215,215}
\definecolor{darkblue}{RGB}{20,70,120}
\hypersetup{colorlinks=true,linkcolor=darkblue,urlcolor=darkblue,citecolor=darkblue}
\pagestyle{fancy}
\fancyhf{}
\lhead{Space Telemetry Ground Station}
\rhead{Technical review}
\cfoot{\thepage}
\setlist[itemize]{leftmargin=1.2em,itemsep=0.2em,topsep=0.2em}
\setlist[enumerate]{leftmargin=1.4em,itemsep=0.2em,topsep=0.2em}
\titleformat{\section}{\Large\bfseries\color{darkblue}}{\thesection}{0.6em}{}
\titleformat{\subsection}{\large\bfseries}{\thesubsection}{0.6em}{}
\lstdefinestyle{cppstyle}{
  language=C++,
  basicstyle=\ttfamily\scriptsize,
  numbers=left,
  numberstyle=\tiny\color{gray},
  stepnumber=1,
  numbersep=6pt,
  backgroundcolor=\color{codebg},
  frame=single,
  rulecolor=\color{codeframe},
  columns=fullflexible,
  keepspaces=true,
  breaklines=true,
  breakatwhitespace=false,
  tabsize=4,
  showstringspaces=false,
  captionpos=b,
  aboveskip=0.8em,
  belowskip=0.8em
}
\lstset{style=cppstyle}
\newcommand{\complexity}[1]{\textbf{Complexity:} #1}
\newcommand{\rationale}[1]{\textbf{Rationale:} #1}
\newcommand{\tradeoff}[1]{\textbf{Trade-off:} #1}
\begin{document}
\begin{titlepage}
\centering
\vspace*{2.5cm}
{\Huge\bfseries Space Telemetry Ground Station - Technical Code Review\par}
\vspace{0.5cm}
{\Large Recruiter-facing technical presentation document\par}
\vspace{1.5cm}
{\large Project: \texttt{cpp-space-telemetry-ground-station}\par}
\vspace{0.4cm}
{\large Date: 3 July 2026\par}
\vfill
\begin{minipage}{0.82\textwidth}
\small

This document explains the project the way I would like a technical recruiter to read it: not as a simple C++ exercise, but as a demonstration of how I design robust software. The code excerpts below were selected to show design logic, algorithmic choices, complexity, error handling, application decomposition and overall development philosophy.
\end{minipage}
\vfill
\end{titlepage}
\tableofcontents
\newpage

\section{Executive summary for recruiters}

\textbf{Space Telemetry Ground Station} is a C++20 Linux-based telemetry ground station simulator. It receives simulated binary telemetry frames over UDP/TCP or from replay files, decodes them, validates integrity with CRC32, exports valid frames to CSV/JSON, and monitors station health through an operational \texttt{DEGRADED} mode.

The goal was not to build a graphical user interface, but to create a credible systems-oriented software base: strict binary protocol, layered architecture, clean CMake, tests, replay, benchmarking, logging, multithreading and documentation.

\subsection{What the project demonstrates}
\begin{itemize}
  \item ability to write modern, readable and testable C++;
  \item understanding of Linux constraints, networking, binary files and robustness;
  \item ability to design a clear and verifiable protocol;
  \item performance reasoning: complexity, throughput and contention effects;
  \item engineering discipline: explicit errors, logs, replay, benchmark, CI and documentation.
\end{itemize}

\subsection{Skill map}
\begin{longtable}{@{}p{0.28\linewidth}p{0.66\linewidth}@{}}
\toprule
Skill & Where it appears in the project \\
\midrule
Modern C++20 & \texttt{std::span}, \texttt{std::variant}, RAII, strong types, library/application split \\
Linux & POSIX sockets, \texttt{poll()}, signals, binary files, CMake build \\
Networking & UDP reception, TCP byte streams, frame extraction \\
Robustness & CRC32, magic bytes, version checks, strict lengths, detailed errors \\
Multithreading & blocking queues, decoder workers, writer thread \\
Performance & reproducible benchmark, throughput and contention analysis \\
Testing & C++ unit tests and CTest integration \\
Documentation & README, technical design, performance report, this document \\
\bottomrule
\end{longtable}

\section{Coding philosophy}

The project follows five principles.

\begin{enumerate}
  \item \textbf{Make boundaries explicit.} Sizes, versions, lengths, formats and thresholds are visible in the code.
  \item \textbf{Keep responsibilities isolated.} The codec does not handle networking; networking does not understand business semantics; the application orchestrates threads and exports.
  \item \textbf{Fail fast, but keep the service alive.} An invalid frame is rejected with a precise error code; the service continues processing subsequent frames.
  \item \textbf{Minimize dependencies when useful.} The project avoids heavy libraries for JSON or networking to show a clear understanding of low-level mechanisms.
  \item \textbf{Measure instead of assuming.} The benchmark and performance report make behavior observable.
\end{enumerate}

\section{Overall architecture}

\subsection{Logical decomposition}
\begin{longtable}{@{}p{0.25\linewidth}p{0.68\linewidth}@{}}
\toprule
Component & Role \\
\midrule
\texttt{stgs\_core} & Protocol, CRC32, codec, replay, logs, health monitoring \\
\texttt{stgs\_network} & UDP/TCP Linux reception through POSIX sockets \\
\texttt{stgs\_ground\_station} & Main application: orchestration, threads, CSV/JSON export \\
\texttt{stgs\_satellite\_simulator} & Frame simulator, loss, corruption, STGF generation \\
\texttt{stgs\_benchmark\_decode} & Reproducible decode throughput measurement \\
\texttt{stgs\_tests} & Unit tests for the functional core \\
\bottomrule
\end{longtable}

\subsection{Data flow}
\begin{verbatim}
Satellite simulator / Replay file
        |
        v
Raw binary frame queue  ---> decoder workers ---> decoded frame queue ---> writer thread
        |                       |                      |                     |
        |                       |                      |                     +--> CSV / JSON
        |                       |                      +--> health monitor
        |                       +--> CRC / protocol validation
        +--> UDP / TCP / STGF
\end{verbatim}

This pipeline was designed like an industrial processing chain: each stage has one clear responsibility, making the code easier to test, maintain and evolve.

\section{Binary protocol design}

\subsection{Excerpt: frame layout}
\begin{lstlisting}[caption={Binary frame layout and protocol constants}]
namespace stgs {

using ByteVector = std::vector<std::uint8_t>;

inline constexpr std::uint32_t FrameMagic = 0x53544753U; // ASCII "STGS"
inline constexpr std::uint8_t ProtocolVersion = 1U;
inline constexpr std::size_t MagicSize = 4U;
inline constexpr std::size_t HeaderSize = 23U;
inline constexpr std::size_t CrcSize = 4U;
inline constexpr std::size_t MinFrameSize = HeaderSize + CrcSize;
inline constexpr std::size_t MaxPayloadSize = 4096U;
inline constexpr std::size_t MaxFrameSize = HeaderSize + MaxPayloadSize + CrcSize;

// Binary frame layout, all integer fields are big-endian:
// MAGIC:u32 | VERSION:u8 | SATELLITE_ID:u16 | TIMESTAMP_MS:u64 |
// TEMPERATURE_C:f32 | BATTERY_PERCENT:u8 | STATUS:u8 | PAYLOAD_LEN:u16 |
// PAYLOAD:bytes | CRC32:u32

enum class Status : std::uint8_t {
    Nominal = 0,
    Warning = 1,
    Critical = 2,
    SafeMode = 3
};

struct TelemetryFrame {
    std::uint8_t version = ProtocolVersion;
    std::uint16_t satelliteId = 0;
    std::uint64_t timestampMs = 0;
\end{lstlisting}


The frame contains a fixed magic value, a version, a satellite identifier, a timestamp, business fields, a payload length and a final CRC. This lets the decoder answer three questions before accepting a frame: \textit{is this the expected protocol?}, \textit{is this the supported version?}, \textit{is the content complete and intact?}

\complexity{Reading fixed header fields is $O(1)$. Reading the payload is $O(p)$ where $p$ is the payload size. The decoded frame uses $O(p)$ memory.}

\rationale{A binary format is stricter than network JSON: it forces sizes, bounds and a single interpretation. This is consistent with telemetry or industrial systems where each byte matters.}

\subsection{Excerpt: portable big-endian serialization}
\begin{lstlisting}[caption={Explicit big-endian serialization}]
inline void appendU8(std::vector<std::uint8_t>& out, std::uint8_t value) {
    out.push_back(value);
}

inline void appendU16BE(std::vector<std::uint8_t>& out, std::uint16_t value) {
    out.push_back(static_cast<std::uint8_t>((value >> 8U) & 0xFFU));
    out.push_back(static_cast<std::uint8_t>(value & 0xFFU));
}

inline void appendU32BE(std::vector<std::uint8_t>& out, std::uint32_t value) {
    out.push_back(static_cast<std::uint8_t>((value >> 24U) & 0xFFU));
    out.push_back(static_cast<std::uint8_t>((value >> 16U) & 0xFFU));
    out.push_back(static_cast<std::uint8_t>((value >> 8U) & 0xFFU));
    out.push_back(static_cast<std::uint8_t>(value & 0xFFU));
}

inline void appendU64BE(std::vector<std::uint8_t>& out, std::uint64_t value) {
    for (int shift = 56; shift >= 0; shift -= 8) {
        out.push_back(static_cast<std::uint8_t>((value >> static_cast<unsigned>(shift)) & 0xFFU));
    }
}

inline void appendFloatBE(std::vector<std::uint8_t>& out, float value) {
    const auto bits = std::bit_cast<std::uint32_t>(value);
    appendU32BE(out, bits);
}

inline std::uint16_t readU16BE(std::span<const std::uint8_t> bytes, std::size_t offset) {
    return static_cast<std::uint16_t>((static_cast<std::uint16_t>(bytes[offset]) << 8U) |
                                      static_cast<std::uint16_t>(bytes[offset + 1U]));
}

\end{lstlisting}


Big-endian, often called network byte order, prevents the format from depending on the host architecture. The code does not directly \texttt{reinterpret\_cast} a C++ structure to bytes, because that would expose the protocol to padding, alignment and endianness issues.

\complexity{Each numeric field is encoded in $O(k)$ where $k$ is the field byte size. Since $k$ is 1, 2, 4 or 8, it is constant-time per field. For a full frame, encoding is dominated by payload copy: $O(p)$.}

\rationale{This is more verbose than raw memory casting, but much more robust and portable.}

\section{Frame encoding and validation}

\subsection{Excerpt: defensive encoding}
\begin{lstlisting}[caption={Frame encoding with precondition checks}]
        if (i > 0U) {
            oss << ' ';
        }
        oss << std::setw(2) << static_cast<int>(payload[i]);
    }
    if (payload.size() > displayed) {
        oss << " ...(" << std::dec << payload.size() << " bytes)";
    }
    return oss.str();
}

std::string toString(const TelemetryFrame& frame) {
    std::ostringstream oss;
    oss << "sat=" << frame.satelliteId
        << " ts_ms=" << frame.timestampMs
        << " temp_c=" << std::fixed << std::setprecision(2) << frame.temperatureC
        << " battery=" << static_cast<int>(frame.batteryPercent) << "%"
        << " status=" << statusToString(frame.status)
        << " payload_len=" << frame.payload.size();
    return oss.str();
}

ByteVector encodeFrame(const TelemetryFrame& frame) {
    if (frame.payload.size() > MaxPayloadSize) {
        throw std::invalid_argument("payload exceeds MaxPayloadSize");
    }
    if (frame.version != ProtocolVersion) {
        throw std::invalid_argument("unsupported protocol version");
    }
    if (frame.batteryPercent > 100U) {
        throw std::invalid_argument("battery percent must be between 0 and 100");
    }

    const auto statusByte = static_cast<std::uint8_t>(frame.status);
    if (statusByte > static_cast<std::uint8_t>(Status::SafeMode)) {
\end{lstlisting}


Encoding starts by validating invariants: maximum payload size, supported version, battery range, and known status. The goal is to never intentionally produce an invalid frame.

\complexity{Encoding reserves memory once, appends fixed fields in $O(1)$, copies the payload in $O(p)$, then computes CRC over the frame in $O(h+p)$. Total complexity is $O(p)$.}

\rationale{Protocol rules are centralized inside the codec. Any future application generating telemetry frames reuses the same safety checks.}

\subsection{Excerpt: strict decoding and explicit errors}
\begin{lstlisting}[caption={Strict decoding and domain errors}]
    }

    ByteVector bytes;
    bytes.reserve(HeaderSize + frame.payload.size() + CrcSize);
    detail::appendU32BE(bytes, FrameMagic);
    detail::appendU8(bytes, frame.version);
    detail::appendU16BE(bytes, frame.satelliteId);
    detail::appendU64BE(bytes, frame.timestampMs);
    detail::appendFloatBE(bytes, frame.temperatureC);
    detail::appendU8(bytes, frame.batteryPercent);
    detail::appendU8(bytes, statusByte);
    detail::appendU16BE(bytes, static_cast<std::uint16_t>(frame.payload.size()));
    bytes.insert(bytes.end(), frame.payload.begin(), frame.payload.end());

    const auto crc = crc32(bytes);
    detail::appendU32BE(bytes, crc);
    return bytes;
}

FrameParseResult decodeFrame(std::span<const std::uint8_t> bytes) {
    if (bytes.size() < MinFrameSize) {
        return makeError(FrameErrorCode::TooShort, "frame is shorter than the minimum header + CRC size");
    }

    const auto magic = detail::readU32BE(bytes, 0U);
    if (magic != FrameMagic) {
        return makeError(FrameErrorCode::BadMagic, "invalid frame magic; expected ASCII STGS");
    }

    const auto version = bytes[4U];
    if (version != ProtocolVersion) {
        return makeError(FrameErrorCode::UnsupportedVersion, "unsupported protocol version");
    }

    const auto payloadLength = detail::readU16BE(bytes, 21U);
    if (payloadLength > MaxPayloadSize) {
        return makeError(FrameErrorCode::PayloadTooLarge, "payload length exceeds configured maximum");
    }

    const auto expectedSize = HeaderSize + static_cast<std::size_t>(payloadLength) + CrcSize;
    if (bytes.size() != expectedSize) {
        std::ostringstream oss;
        oss << "frame length mismatch: expected " << expectedSize << " bytes, got " << bytes.size();
        return makeError(FrameErrorCode::LengthMismatch, oss.str());
    }

    const auto expectedCrc = detail::readU32BE(bytes, bytes.size() - CrcSize);
    const auto calculatedCrc = crc32(bytes.subspan(0U, bytes.size() - CrcSize));
    if (expectedCrc != calculatedCrc) {
        std::ostringstream oss;
        oss << "CRC mismatch: expected 0x" << std::hex << std::setw(8) << std::setfill('0') << expectedCrc
            << ", calculated 0x" << std::setw(8) << calculatedCrc;
        return makeError(FrameErrorCode::BadCrc, oss.str());
    }

    const auto battery = bytes[19U];
    if (battery > 100U) {
        return makeError(FrameErrorCode::InvalidBattery, "battery percent is outside the 0..100 range");
    }

    const auto statusByte = bytes[20U];
    if (statusByte > static_cast<std::uint8_t>(Status::SafeMode)) {
        return makeError(FrameErrorCode::InvalidStatus, "status field contains an unknown value");
    }

\end{lstlisting}


The decoder is written as a sequence of safety gates. It rejects frames that are too short, have a bad magic value, unsupported version, oversized payload, inconsistent length, invalid CRC or invalid business fields.

The result type is \texttt{std::variant<TelemetryFrame, FrameError>}. This avoids mixing exceptions with normal data processing. An invalid frame is expected in a networked system; it should not necessarily crash the service.

\complexity{Decoding is $O(p)$ because CRC must scan all non-CRC bytes. Header checks are $O(1)$. Payload copy is $O(p)$.}

\tradeoff{Using a \texttt{variant} forces callers to handle both success and failure explicitly. This is clearer than returning null or throwing on every corrupted frame.}

\section{CRC32 integrity check}

\subsection{Excerpt: compile-time CRC table}
\begin{lstlisting}[caption={Table-driven CRC32}]
    std::array<std::uint32_t, 256> table{};
    for (std::uint32_t i = 0; i < table.size(); ++i) {
        std::uint32_t crc = i;
        for (int bit = 0; bit < 8; ++bit) {
            if ((crc & 1U) != 0U) {
                crc = (crc >> 1U) ^ 0xEDB88320U;
            } else {
                crc >>= 1U;
            }
        }
        table[i] = crc;
    }
    return table;
}

constexpr auto CrcTable = makeTable();

} // namespace

std::uint32_t crc32(std::span<const std::uint8_t> bytes) noexcept {
    std::uint32_t crc = 0xFFFFFFFFU;
    for (const auto byte : bytes) {
        const auto index = static_cast<std::uint8_t>((crc ^ byte) & 0xFFU);
        crc = (crc >> 8U) ^ CrcTable[index];
    }
\end{lstlisting}


CRC32 uses a 256-entry \texttt{constexpr} table. The table avoids recomputing eight bit-level steps for each input byte at runtime. The \texttt{0xEDB88320} polynomial corresponds to a common CRC32 variant.

\complexity{The computation is $O(n)$ for $n$ bytes. Extra memory is $O(1)$: a fixed 256-entry table, about 1 KiB. Without the table, the complexity would still be $O(n)$ but with a higher constant factor.}

\rationale{CRC32 is not cryptographic. It detects corruption, truncation or accidental modification. Authentication would require a MAC or signature.}

\section{Frame extraction from TCP streams}

\subsection{Excerpt: reassembling complete frames}
\begin{lstlisting}[caption={Complete frame extraction from a TCP stream}]
        return "LengthMismatch";
    case FrameErrorCode::BadCrc:
        return "BadCrc";
    case FrameErrorCode::InvalidBattery:
        return "InvalidBattery";
    case FrameErrorCode::InvalidStatus:
        return "InvalidStatus";
    }
    return "Unknown";
}

std::vector<ByteVector> StreamFrameExtractor::feed(std::span<const std::uint8_t> bytes) {
    buffer_.insert(buffer_.end(), bytes.begin(), bytes.end());
    std::vector<ByteVector> frames;
    const auto magic = magicBytes();

    while (true) {
        const auto it = std::search(buffer_.begin(), buffer_.end(), magic.begin(), magic.end());
        if (it == buffer_.end()) {
            if (buffer_.size() > MagicSize - 1U) {
                buffer_.erase(buffer_.begin(), buffer_.end() - static_cast<std::ptrdiff_t>(MagicSize - 1U));
            }
            return frames;
        }

        if (it != buffer_.begin()) {
            buffer_.erase(buffer_.begin(), it);
        }

        if (buffer_.size() < HeaderSize) {
            return frames;
        }

        const auto payloadLength = detail::readU16BE(buffer_, 21U);
        if (payloadLength > MaxPayloadSize) {
            buffer_.erase(buffer_.begin());
            continue;
        }

        const auto frameSize = HeaderSize + static_cast<std::size_t>(payloadLength) + CrcSize;
        if (buffer_.size() < frameSize) {
            return frames;
        }

        frames.emplace_back(buffer_.begin(), buffer_.begin() + static_cast<std::ptrdiff_t>(frameSize));
        buffer_.erase(buffer_.begin(), buffer_.begin() + static_cast<std::ptrdiff_t>(frameSize));
    }
\end{lstlisting}


UDP preserves datagram boundaries. TCP provides a byte stream without message boundaries. The \texttt{StreamFrameExtractor} therefore keeps an internal buffer, searches for the \texttt{STGS} magic value, waits for a complete header, reads the payload length, and returns only complete frames.

\complexity{Each feed inserts $n$ received bytes in $O(n)$. Magic search with \texttt{std::search} is practically close to $O(n)$ because the pattern has 4 bytes. The theoretical worst case is $O(n \cdot m)$ with $m=4$, effectively $O(n)$.}

\tradeoff{The implementation favors readability. For very high throughput, a ring buffer could reduce copies and memory moves.}

\section{Multithreading: producer/consumer pipeline}

\subsection{Excerpt: generic blocking queue}
\begin{lstlisting}[caption={Thread-safe queue with explicit closure}]
class BlockingQueue {
public:
    BlockingQueue() = default;
    BlockingQueue(const BlockingQueue&) = delete;
    BlockingQueue& operator=(const BlockingQueue&) = delete;

    bool push(T value) {
        {
            std::lock_guard<std::mutex> lock(mutex_);
            if (closed_) {
                return false;
            }
            queue_.push(std::move(value));
        }
        cv_.notify_one();
        return true;
    }

    std::optional<T> pop() {
        std::unique_lock<std::mutex> lock(mutex_);
        cv_.wait(lock, [this] { return closed_ || !queue_.empty(); });
        if (queue_.empty()) {
            return std::nullopt;
        }
        T value = std::move(queue_.front());
        queue_.pop();
        return value;
    }

    void close() {
        {
            std::lock_guard<std::mutex> lock(mutex_);
            closed_ = true;
        }
        cv_.notify_all();
    }

    [[nodiscard]] bool closed() const {
        std::lock_guard<std::mutex> lock(mutex_);
        return closed_;
    }

    [[nodiscard]] std::size_t size() const {
        std::lock_guard<std::mutex> lock(mutex_);
        return queue_.size();
    }

private:
    mutable std::mutex mutex_;
\end{lstlisting}


The blocking queue decouples reception, decoding and writing. \texttt{std::condition\_variable} prevents busy-waiting: workers sleep until work is available or the queue is closed.

\complexity{\texttt{push()} and \texttt{pop()} are amortized $O(1)$. Memory is $O(q)$ where $q$ is the number of queued elements. Latency depends mostly on mutex contention and copy/move costs.}

\rationale{\texttt{close()} is important: consumers can exit cleanly without sending artificial sentinel values through the queue.}

\subsection{Excerpt: worker orchestration}
\begin{lstlisting}[caption={Decoder workers and writer thread}]
        std::thread writerThread([&] {
            bool firstJsonFrame = true;
            while (auto frame = decodedQueue.pop()) {
                if (outputFormat == OutputFormat::Json) {
                    writeFrameJson(out, *frame, firstJsonFrame);
                    firstJsonFrame = false;
                } else {
                    writeFrameCsv(out, *frame);
                }
                ++written;
            }
            if (outputFormat == OutputFormat::Json) {
                writeJsonFooter(out);
            }
            out.flush();
        });

        std::vector<std::thread> decoderThreads;
        decoderThreads.reserve(opts.decoderThreads);
        for (std::size_t i = 0; i < opts.decoderThreads; ++i) {
            decoderThreads.emplace_back([&, i] {
                while (auto bytes = rawQueue.pop()) {
                    auto result = stgs::decodeFrame(*bytes);
                    if (std::holds_alternative<stgs::TelemetryFrame>(result)) {
                        auto frame = std::get<stgs::TelemetryFrame>(std::move(result));
                        ++decoded;
                        if (auto transition = health.recordDecoded(frame); transition.has_value()) {
                            logTransition(logger, *transition);
                        }
                        decodedQueue.push(std::move(frame));
                    } else {
                        const auto& err = std::get<stgs::FrameError>(result);
                        ++rejected;
                        if (auto transition = health.recordRejected(); transition.has_value()) {
                            logTransition(logger, *transition);
                        }
                        logger.warning("decoder " + std::to_string(i) + " rejected frame: " +
                                       stgs::errorCodeToString(err.code) + " - " + err.message);
                    }
                }
            });
        }
\end{lstlisting}


The pipeline separates CPU cost from I/O cost. Workers consume raw binary frames, call \texttt{decodeFrame()}, and push valid frames into the decoded queue. The writer serializes output to CSV or JSON.

\complexity{For $N$ frames with average payload size $p$, total work is $O(N \cdot p)$. With $T$ workers, ideal time tends toward $O((N \cdot p)/T)$, but only until queue contention, allocations and I/O dominate.}

\tradeoff{The benchmark shows that adding too many threads can reduce throughput. Multithreading must be measured, not assumed.}

\section{Degraded mode: operational monitoring}

\subsection{Excerpt: sliding window and hysteresis}
\begin{lstlisting}[caption={NOMINAL / DEGRADED state transitions}]
        throw std::invalid_argument("recovery rejection rate must be between 0 and 1");
    }
    if (config_.recoveryRejectionRate > config_.degradedRejectionRate) {
        throw std::invalid_argument("recovery rejection rate cannot exceed degraded rejection rate");
    }
}

std::optional<StationStateTransition> StationHealthMonitor::recordDecoded(const TelemetryFrame& frame) {
    return recordSample(Sample{false, isCriticalTelemetry(frame)});
}

std::optional<StationStateTransition> StationHealthMonitor::recordRejected() {
    return recordSample(Sample{true, false});
}

std::optional<StationStateTransition> StationHealthMonitor::recordSample(Sample sample) {
    std::lock_guard<std::mutex> lock(mutex_);
    if (!config_.enabled) {
        return std::nullopt;
    }

    samples_.push_back(sample);
    while (samples_.size() > config_.windowSize) {
        samples_.pop_front();
    }

    const auto snapshot = makeSnapshotLocked();
    if (snapshot.samples < config_.minSamples) {
        return std::nullopt;
    }

    StationState target = state_;
    if (state_ == StationState::Nominal) {
        if (snapshot.rejectionRate >= config_.degradedRejectionRate ||
            snapshot.critical >= config_.criticalFramesForDegraded) {
            target = StationState::Degraded;
        }
    } else {
        if (snapshot.rejectionRate <= config_.recoveryRejectionRate && snapshot.critical == 0U) {
            target = StationState::Nominal;
        }
    }

    if (target == state_) {
        return std::nullopt;
    }

    StationStateTransition transition;
    transition.from = state_;
    transition.to = target;
    transition.snapshot = snapshot;
    transition.reason = transitionReasonLocked(snapshot, target);
    state_ = target;
    transition.snapshot.state = state_;
    return transition;
}

\end{lstlisting}


\texttt{DEGRADED} is the station state, not just a telemetry frame status. It is triggered when rejection rate exceeds a threshold or when several critical frames appear in the monitoring window. Recovery uses a lower threshold, creating hysteresis and avoiding rapid oscillations.

\complexity{The current implementation computes a snapshot by scanning the window: $O(w)$ per sample, with $w$ the window size. Default $w=100$, so cost is bounded and small. A future optimization could keep incremental counters and reach $O(1)$ per sample.}

\rationale{The implementation favors readability and auditability. In critical environments, a bounded and understandable algorithm can be preferable to premature micro-optimization.}

\section{Binary STGF replay}

\subsection{Excerpt: defensive replay format}
\begin{lstlisting}[caption={Reading and writing STGF replay files}]
void FrameFileWriter::writeFrame(std::span<const std::uint8_t> frameBytes) {
    if (frameBytes.size() > MaxFrameSize) {
        throw std::runtime_error("refusing to write frame larger than MaxFrameSize");
    }
    ByteVector length;
    detail::appendU32BE(length, static_cast<std::uint32_t>(frameBytes.size()));
    out_.write(reinterpret_cast<const char*>(length.data()), static_cast<std::streamsize>(length.size()));
    out_.write(reinterpret_cast<const char*>(frameBytes.data()), static_cast<std::streamsize>(frameBytes.size()));
    if (!out_) {
        throw std::runtime_error("failed while writing replay frame");
    }
}

FrameFileReader::FrameFileReader(const std::filesystem::path& path) {
    in_.open(path, std::ios::binary | std::ios::in);
    if (!in_) {
        throw std::runtime_error("failed to open replay input file: " + path.string());
    }

    std::array<std::uint8_t, FileHeader.size()> header{};
    readExact(in_, header.data(), header.size());
    if (header != FileHeader) {
        throw std::runtime_error("invalid or unsupported STGF replay file header");
    }
}

std::optional<ByteVector> FrameFileReader::readNext() {
    std::array<std::uint8_t, 4> lenBytes{};
    in_.read(reinterpret_cast<char*>(lenBytes.data()), static_cast<std::streamsize>(lenBytes.size()));
    const auto bytesRead = in_.gcount();
    if (bytesRead == 0 && in_.eof()) {
        return std::nullopt;
    }
    if (bytesRead != static_cast<std::streamsize>(lenBytes.size())) {
        throw std::runtime_error("truncated frame length in replay file");
    }

    const auto length = detail::readU32BE(lenBytes, 0U);
    if (length < MinFrameSize || length > MaxFrameSize) {
        throw std::runtime_error("invalid replay frame length");
    }

    ByteVector frame(length);
    readExact(in_, frame.data(), frame.size());
    return frame;
\end{lstlisting}


Replay makes incidents reproducible. A \texttt{STGF} file starts with a fixed header, then stores each frame prefixed by its length. Reading validates the header, minimum size, maximum size and truncated end-of-file cases.

\complexity{Reading and writing are $O(n)$ for $n$ stored bytes. Each frame adds a constant 4-byte length overhead.}

\rationale{Replay is critical for industrial debugging: instead of saying that a network issue occurred, captured frames can be replayed exactly.}

\section{CSV/JSON export without external dependency}

\subsection{Excerpt: JSON escaping and output writing}
\begin{lstlisting}[caption={CSV/JSON export}]
std::string jsonEscape(const std::string& value) {
    std::ostringstream out;
    for (unsigned char ch : value) {
        switch (ch) {
        case '"':
            out << "\\\"";
            break;
        case '\\':
            out << "\\\\";
            break;
        case '\b':
            out << "\\b";
            break;
        case '\f':
            out << "\\f";
            break;
        case '\n':
            out << "\\n";
            break;
        case '\r':
            out << "\\r";
            break;
        case '\t':
            out << "\\t";
            break;
        default:
            if (ch < 0x20U) {
                out << "\\u" << std::hex << std::setw(4) << std::setfill('0') << static_cast<int>(ch);
            } else {
                out << static_cast<char>(ch);
            }
            break;
        }
    }
    return out.str();
}

void writeCsvHeader(std::ofstream& out) {
    out << "timestamp_ms,satellite_id,temperature_c,battery_percent,status,payload_len,payload_hex\n";
}

void writeFrameCsv(std::ofstream& out, const stgs::TelemetryFrame& frame) {
    out << frame.timestampMs << ','
        << frame.satelliteId << ','
        << frame.temperatureC << ','
        << static_cast<int>(frame.batteryPercent) << ','
        << stgs::statusToString(frame.status) << ','
        << frame.payload.size() << ','
        << csvEscape(stgs::payloadToHex(frame.payload, frame.payload.size())) << '\n';
}

void writeJsonHeader(std::ofstream& out) {
    out << "[\n";
}

void writeFrameJson(std::ofstream& out, const stgs::TelemetryFrame& frame, bool first) {
    if (!first) {
        out << ",\n";
    }
    out << "  {"
        << "\"timestamp_ms\":" << frame.timestampMs << ','
        << "\"satellite_id\":" << frame.satelliteId << ','
        << "\"temperature_c\":" << std::setprecision(6) << frame.temperatureC << ','
        << "\"battery_percent\":" << static_cast<int>(frame.batteryPercent) << ','
        << "\"status\":\"" << jsonEscape(stgs::statusToString(frame.status)) << "\"," 
        << "\"payload_len\":" << frame.payload.size() << ','
        << "\"payload_hex\":\"" << jsonEscape(stgs::payloadToHex(frame.payload, frame.payload.size())) << "\""
        << '}';
}

void writeJsonFooter(std::ofstream& out) {
    out << "\n]\n";
\end{lstlisting}


JSON output is implemented without an external library to keep the project portable. The sensitive part is string escaping: quotes, backslashes and control characters. CSV is also escaped to preserve column structure.

\complexity{Escaping is $O(s)$ where $s$ is string length. Total export is $O(N \cdot p)$ because payload-to-hex conversion scans each payload.}

\tradeoff{For a larger production application, a well-tested JSON library would be preferable. Here the manual implementation is educational and limits dependencies.}

\section{Linux network reception}

\subsection{Excerpt: UDP loop with poll}
\begin{lstlisting}[caption={Linux UDP loop using poll}]

    int yes = 1;
    if (::setsockopt(fd, SOL_SOCKET, SO_REUSEADDR, &yes, sizeof(yes)) < 0) {
        ::close(fd);
        throw std::runtime_error("setsockopt(SO_REUSEADDR) failed: " + std::string(std::strerror(errno)));
    }

    const auto addr = makeAddress(config_.bindAddress, config_.port);
    if (::bind(fd, reinterpret_cast<const sockaddr*>(&addr), sizeof(addr)) < 0) {
        ::close(fd);
        throw std::runtime_error("bind() failed on " + config_.bindAddress + ":" +
                                 std::to_string(config_.port) + ": " + std::strerror(errno));
    }

    return fd;
}

void NetworkServer::runUdp(FrameCallback& callback, const std::atomic_bool& running) {
    serverFd_ = createBoundSocket(SOCK_DGRAM);
    logger_.info("UDP receiver listening on " + config_.bindAddress + ":" + std::to_string(config_.port));

    ByteVector buffer(MaxFrameSize);
    while (running.load() && !stopRequested_.load()) {
        pollfd pfd{serverFd_, POLLIN, 0};
        const int rc = ::poll(&pfd, 1, config_.pollTimeoutMs);
        if (rc < 0) {
            if (errno == EINTR) {
                continue;
            }
            throw std::runtime_error("poll() failed: " + std::string(std::strerror(errno)));
        }
        if (rc == 0 || (pfd.revents & POLLIN) == 0) {
            continue;
        }

        sockaddr_in src{};
        socklen_t srcLen = sizeof(src);
        const ssize_t n = ::recvfrom(serverFd_, buffer.data(), buffer.size(), 0,
                                     reinterpret_cast<sockaddr*>(&src), &srcLen);
        if (n < 0) {
\end{lstlisting}


UDP reception uses POSIX sockets and \texttt{poll()} with a timeout. This allows the loop to wait for data without blocking forever, while still checking the stop signal. Transient errors such as \texttt{EINTR} are handled without stopping the server.

\complexity{For UDP, each datagram is processed in $O(f)$ where $f$ is frame size. For TCP, polling is $O(c)$ per loop with $c$ connected clients.}

\tradeoff{\texttt{poll()} is simple and portable on Linux/Unix. With thousands of clients, \texttt{epoll()} would be more appropriate.}

\section{Benchmark and performance analysis}

\subsection{Excerpt: decode pipeline benchmark}
\begin{lstlisting}[caption={Decode pipeline benchmark}]
int main(int argc, char** argv) {
    try {
        const auto opts = parseArgs(argc, argv);
        std::mt19937 rng(opts.seed);
        std::vector<stgs::ByteVector> encoded;
        encoded.reserve(opts.frames);
        for (std::size_t i = 0; i < opts.frames; ++i) {
            encoded.push_back(stgs::encodeFrame(makeFrame(opts, rng, i)));
        }

        stgs::BlockingQueue<stgs::ByteVector> queue;
        std::atomic_size_t decoded{0};
        std::atomic_size_t rejected{0};

        std::vector<std::thread> workers;
        workers.reserve(opts.decoderThreads);
        const auto start = std::chrono::steady_clock::now();
        for (std::size_t i = 0; i < opts.decoderThreads; ++i) {
            workers.emplace_back([&] {
                while (auto bytes = queue.pop()) {
                    auto result = stgs::decodeFrame(*bytes);
                    if (std::holds_alternative<stgs::TelemetryFrame>(result)) {
                        ++decoded;
                    } else {
                        ++rejected;
                    }
                }
            });
        }

        for (const auto& frame : encoded) {
            queue.push(frame);
        }
        queue.close();
        for (auto& worker : workers) {
            worker.join();
        }
        const auto end = std::chrono::steady_clock::now();
        const auto seconds = std::chrono::duration<double>(end - start).count();
        const auto fps = static_cast<double>(decoded.load() + rejected.load()) / seconds;

        std::cout << "benchmark=decode_pipeline"
                  << " frames=" << opts.frames
                  << " payload_size=" << opts.payloadSize
                  << " decoder_threads=" << opts.decoderThreads
                  << " decoded=" << decoded.load()
                  << " rejected=" << rejected.load()
                  << " duration_seconds=" << seconds
                  << " frames_per_second=" << fps << '\n';

\end{lstlisting}


The benchmark generates valid frames, pushes them through the queue, starts decoder workers and measures throughput. It also checks that all expected frames were decoded and none were rejected.

\complexity{The measured workload is $O(N \cdot p)$ with $N$ frames and payload size $p$. Throughput comparison by payload size shows the impact of CRC and payload copies.}

\subsection{Measured results in the generation environment}

\begin{longtable}{@{}rrrr@{}}
\toprule
Frames & Payload & Decoder threads & Measured throughput \\
\midrule
200000 & 64 B & 1 & 1091780 frames/s \\
200000 & 64 B & 2 & 1133140 frames/s \\
200000 & 64 B & 4 & 135736 frames/s \\
100000 & 16 B & 2 & 1035120 frames/s \\
100000 & 64 B & 2 & 1238090 frames/s \\
100000 & 256 B & 2 & 438356 frames/s \\
100000 & 1024 B & 2 & 305913 frames/s \\
100000 & 4096 B & 2 & 93070 frames/s \\
\bottomrule
\end{longtable}


These figures are not absolute performance guarantees; they depend on machine, compiler and system load. They are useful to analyze trends.

\begin{itemize}
  \item A 4096-byte payload strongly reduces throughput, as expected for linear CRC and copy work.
  \item Moving from 1 to 2 threads is only slightly positive on the 64-byte case.
  \item Moving to 4 threads is significantly worse in this run: queue contention, scheduling and copies probably dominate.
\end{itemize}

The key engineering conclusion is: \textbf{multithreading is not automatically an optimization}. It must be measured.

\section{Tests and quality}

\subsection{Excerpt: protocol and error tests}
\begin{lstlisting}[caption={Tests for CRC, round-trip and errors}]
    auto parsed = stgs::decodeFrame(bytes);
    ASSERT_TRUE(std::holds_alternative<stgs::FrameError>(parsed));
    ASSERT_EQ(std::get<stgs::FrameError>(parsed).code, stgs::FrameErrorCode::BadCrc);
}

void testLengthMismatch() {
    auto bytes = stgs::encodeFrame(sampleFrame());
    bytes.pop_back();
    auto parsed = stgs::decodeFrame(bytes);
    ASSERT_TRUE(std::holds_alternative<stgs::FrameError>(parsed));
    ASSERT_EQ(std::get<stgs::FrameError>(parsed).code, stgs::FrameErrorCode::LengthMismatch);
}

void testStreamExtractorSplitAndNoise() {
    const auto frameA = stgs::encodeFrame(sampleFrame());
    auto frameBValue = sampleFrame();
    frameBValue.satelliteId = 99;
    frameBValue.payload = {1, 2, 3};
    const auto frameB = stgs::encodeFrame(frameBValue);

    stgs::StreamFrameExtractor extractor;
    stgs::ByteVector chunk1 = {0x00, 0x11, 0x22};
    chunk1.insert(chunk1.end(), frameA.begin(), frameA.begin() + 10);
\end{lstlisting}


Tests cover essential properties: known CRC32 vector, encode/decode round-trip, bad magic, bad CRC, length mismatch, TCP extraction, replay, degraded mode and blocking queue.

\complexity{Codec tests are $O(p)$ per frame, matching the tested implementation. The goal is to validate invariants rather than exhaustively prove every possible input.}

\rationale{Tests document expected protocol guarantees: corrupted frames must be rejected, valid frames must be recovered faithfully.}

\subsection{Excerpt: CMake and CTest}
\begin{lstlisting}[caption={CMake structure: libraries, executables, tests, benchmarks}]
cmake_minimum_required(VERSION 3.16)

project(cpp_space_telemetry_ground_station
    DESCRIPTION "C++20 Linux telemetry ground station simulator"
    LANGUAGES CXX)

set(CMAKE_CXX_STANDARD 20)
set(CMAKE_CXX_STANDARD_REQUIRED ON)
set(CMAKE_CXX_EXTENSIONS OFF)

option(STGS_BUILD_TESTS "Build unit tests" ON)
option(STGS_BUILD_BENCHMARKS "Build benchmark tools" ON)

find_package(Threads REQUIRED)

set(STGS_WARNINGS -Wall -Wextra -Wpedantic)
if(CMAKE_CXX_COMPILER_ID STREQUAL "GNU")
    list(APPEND STGS_WARNINGS -Wno-free-nonheap-object)
endif()

add_library(stgs_core
    src/Crc32.cpp
    src/FrameCodec.cpp
    src/Logger.cpp
    src/Replay.cpp
    src/StationHealth.cpp)

target_include_directories(stgs_core
    PUBLIC
        ${CMAKE_CURRENT_SOURCE_DIR}/include)

target_compile_options(stgs_core PRIVATE ${STGS_WARNINGS})

add_library(stgs_network
    src/NetworkServer.cpp)

target_link_libraries(stgs_network
    PUBLIC
        stgs_core
        Threads::Threads)

target_compile_options(stgs_network PRIVATE ${STGS_WARNINGS})

add_executable(stgs_ground_station
    apps/ground_station.cpp)

target_link_libraries(stgs_ground_station
    PRIVATE
        stgs_network
        stgs_core
        Threads::Threads)
\end{lstlisting}


The project is split into libraries and executables. This avoids linking the full application just to test the core. It also supports future reuse of \texttt{stgs\_core} by another tool.

\section{Known limits and possible evolutions}

\begin{longtable}{@{}p{0.32\linewidth}p{0.60\linewidth}@{}}
\toprule
Current limit & Industrial evolution \\
\midrule
Single-mutex queue & Lock-free or optimized MPMC queue if benchmarks require it \\
$O(w)$ degraded snapshot & Incremental counters for $O(1)$ per sample \\
Manual JSON & Robust JSON library if the output format grows \\
\texttt{poll()} for TCP & \texttt{epoll()} for many clients \\
CRC32 is not cryptographic & HMAC/signature if authenticity is required \\
No database persistence & PostgreSQL or time-series storage \\
No Linux packaging & Install target, systemd service, Debian package \\
\bottomrule
\end{longtable}

\section{What a recruiter should understand}

The project shows a developer mindset that goes beyond making a program work. The visible choices are those of software meant to be explained, tested, replayed, measured and maintained.

\begin{itemize}
  \item I accept useful complexity: binary protocol, CRC, replay, multithreading.
  \item I avoid gratuitous complexity: no heavy framework hiding the mechanisms.
  \item I think about errors: corruption, truncation, invalid versions, clean thread shutdown.
  \item I think about operations: logs, degraded mode, replay, benchmark.
  \item I think about technical recruitment: the code must be understandable by someone who has never met me.
\end{itemize}

\section{Conclusion}

This telemetry ground station is an industrial-oriented C++ portfolio project. It does not claim to be certified space software, but it demonstrates a concrete base: C++20, Linux, networking, binary protocol, CRC32, multithreading, tests, benchmark and documentation. Its main value is to show my coding philosophy: explicit, measurable, robust, documented and maintainability-oriented.

\end{document}
